\rSec3[macro.general]{Macro continuations}

\begin{codeblock}
namespace std {
  template<class T>
  class @\libglobal{gadget}@ {
  public:
    int size() const;
  };
}
\end{codeblock}

\pnum
\placeholder{noex} is a placeholder in leading prose.

\indexlibrarymember{size}{gadget}%
\begin{itemdecl}
int size() const;
\end{itemdecl}

\begin{itemdescr}
\pnum
\expects
\tcode{T} meets the \oldconcept{Destructible} and \oldconcept{CopyConstructible} requirements.

\pnum
\returns
\tcode{0}.

\pnum
\remarks
The following descriptions apply:
\begin{itemize}
\item Meets \oldconcept{CopyConstructible} requirements.
\end{itemize}
\placeholder{after} follows the list.
\begin{libtab2}{Flat table \placeholder{caption}}{macro.flat}{lp{4.5in}}{Type \cv{} qualifier}{Meaning}

A \placeholder{key} &
Meets \oldconcept{Destructible} requirements. \\
\end{libtab2}
\begin{lib2dtab2}{Grid table}{macro.grid}
{A}
{B}

\rowhdr{Type \cv{} qualifier} &
First \placeholder{value} &
Second \placeholder{value} \\
\end{lib2dtab2}
\placeholder{resume} follows the tables.
\begin{note}
The type meets \oldconcept{Destructible} requirements.
\end{note}
\begin{example}
A \cv{} qualifier is written literally.
\end{example}
\end{itemdescr}
